% =====================================================================
%  Цель и задание лабораторной работы.
% =====================================================================

\textbf{Цель:} освоить на практике основные принципы машинного перевода
документов.

\vspace{1em}

\textbf{Задание (вариант 11)}: разработать систему автоматического
машинного перевода текстов со следующей комбинацией признаков варианта:
\begin{itemize}
  \item \textbf{Направление перевода} -- англо-французский;
  \item \textbf{Предметная область} -- научные статьи по медицине,
  критика предметов изобразительного искусства.
\end{itemize}

\vspace{1em}

Система должна:
\begin{itemize}
  \item на входе принимать естественно-языковой текст на входном языке,
  подлежащий машинному переводу;
  \item подсчитывать количество слов во входном тексте, количество
  слов в переводе и долю слов, для которых перевод найден в словаре, а
  также определять грамматическую информацию (теги частей речи -- с
  использованием функциональности ЛР3 весеннего семестра, POS-теги,
  выводятся как есть, без перевода на русский язык);
  \item на выходе выдавать:
  \begin{enumerate}
    \item перевод входного текста на выходной язык;
    \item упорядоченный по частоте встречаемости в тексте список слов
    и их переводов на выходной язык с грамматической информацией
    (\emph{вкладка 1}, с использованием функциональности ЛР1
    весеннего семестра -- частотный список слов);
    \item построенное дерево синтаксического разбора выбранного
    предложения (\emph{вкладка 2});
  \end{enumerate}
  \item обеспечивать наличие утилиты автоматического
  пополнения/корректировки полученного словаря (таблицы БД);
  \item обеспечивать сохранение и распечатку результатов перевода и
  упорядоченных по частоте встречаемости в тексте списков слов и их
  переводов на выходной язык с грамматической информацией в файл
  формата txt кодировки Unicode;
  \item иметь простой и доступный пользовательский интерфейс с понятным
  набором инструментов.
\end{itemize}

Среди трёх основных типов систем машинного перевода -- систем прямого
перевода, систем непрямого перевода (с трансфером или
языком-посредником) и систем, основанных на знаниях -- по варианту для
реализации выбрана \textbf{система прямого (пословного) перевода}:
исходный текст токенизируется, каждое слово ищется в двуязычном словаре
по паре (лемма, часть речи) и заменяется на переводной эквивалент;
переводная модель как таковая не используется, учитывается лишь
локальный контекст (грамматическая форма слова через POS-тег). Сверх
обязательного по варианту дополнительно реализованы ещё два типа --
непрямой перевод с трансфером и нейросетевой перевод предобученной
моделью -- для практического сравнения всех трёх на одной тестовой
коллекции (см. раздел «Трансферный и нейросетевой методы» далее).
